\section{Smoothing, energy and moments}\label{wt:sec:round2}

\subsection{The heat trace as a Gaussian smoothing of the prime error}

Let $D(u)=(\psi_{\mathrm C}(e^u)-e^u)e^{-u/2}$ be the normalised prime error on the log scale, so $D(0)=-1$.

\begin{proposition}\label{wt:prop:smooth}
For every $t>0$ (without RH),
\begin{equation}\label{wt:eq:smooth}
W(t)=\Sigma_\infty(t)+\int_0^\infty D(u)\,k_t(u)\,du,\qquad k_t=\partial_ug_t-\tfrac12g_t,
\end{equation}
with
\[
\Sigma_\infty(t)=\int_0^\infty g_t(u)e^{-u/2}du-g_t(0)\big(1+\tfrac12\log\pi\big)+\frac1{4\pi}\int e^{-tr^2}\Re\psi\big(\tfrac14+\tfrac{ir}2\big)dr.
\]
\end{proposition}

\begin{proof}
Write $d\psi_{\mathrm C}(e^u)=e^u\,du+d(D\,e^{u/2})$ in \eqref{wt:eq:P} and integrate the second part by parts: $\mathcal P=\int_0^\infty g_te^{u/2}du-g_t(0)D(0)-\int_0^\infty D\,(\partial_ug_t-\frac12g_t)\,du$, the upper boundary vanishing by Gaussian decay and $\psi_{\mathrm C}(y)\ll y$. Since $\int_0^\infty g_t(e^{u/2}+e^{-u/2})\,du=\int_{\R}g_te^{u/2}du=e^{t/4}$, subtracting $\mathcal P$ from \eqref{wt:eq:A} gives the claim.
\end{proof}

The source verified \eqref{wt:eq:smooth} against the zeros to 18 digits at $t=0.01,0.05,0.1,0.2$, where $\Sigma_\infty$ is of order one and negative: $-3.20$, $-1.77$, $-1.33$, $-0.99$ (unchecked here). So the Gaussian derivative smoothing of $D$ at scale $\sqrt t$ must cancel an explicit archimedean quantity to exponential accuracy. By the explicit formula $D(u)\approx-\sum_\rho e^{(\rho-1/2)u}/\rho$, and the bilateral Gaussian transform $\int_{\R}g_t(u)e^{qu}du=e^{tq^2}$, for $q=\rho-\frac12=i\gamma$, is the factor $e^{-t\gamma^2}$ of the heat trace, complex off the critical line; this explains the heat factor, although the half line integral in \eqref{wt:eq:smooth} also carries a derivative, boundary and archimedean terms, so it is not a pure Fourier multiplier on the zero series. Under RH $W(t)\sim m_{\gamma_1}e^{-\gamma_1^2t}$ as $t\to\infty$, and an off line zero, with growth $e^{\delta u}$ in $D$, would spoil complete monotonicity. Integrating $W=\mathcal A-\mathcal P$ term by term is valid on a finite interval $(0,X]$; it must not be read as separate integration over $(0,\infty)$, since $\int_0^\infty g_t(u)\,dt=\infty$ and $\mathcal A$ contains $e^{t/4}$, and the convergent full integral comes from the cancellation in $\mathcal A-\mathcal P$.

\subsection{Variance, energy and RH}

\begin{proposition}[Energy form of RH]\label{wt:prop:energy}
RH holds if and only if $E_\epsilon=\int_0^\infty D(u)^2e^{-2\epsilon u}\,du<\infty$ for every $\epsilon>0$. Under RH,
\[
\lim_{U\to\infty}\frac1U\int_0^UD(u)^2\,du=2\sum_{\gamma>0}\frac{m_\gamma^2}{\frac14+\gamma^2},
\]
the sum over distinct positive ordinates, with multiplicities $m_\gamma$. This equals $\sum_\rho|\rho|^{-2}=2+\gamma_E-\log4\pi=0.0461914179$ if all zeros are simple, and equals $2+\gamma_E-\log4\pi+2\sum_{\gamma>0}m_\gamma(m_\gamma-1)/(\frac14+\gamma^2)$ in general.
\end{proposition}

\begin{proof}
For $\Re s>1$, $\int_0^\infty D(u)e^{-zu}du=\int_1^\infty(\psi_{\mathrm C}(y)-y)y^{-s-1}dy=-\frac{\zeta'(s)}{s\zeta(s)}-\frac1{s-1}$ with $s=\frac12+z$, by the substitution $y=e^u$ and the Dirichlet series of $-\zeta'/\zeta$; the singularity at $s=1$ cancels, and at a zero $\rho$ of multiplicity $m_\rho$ the residue is $-m_\rho/\rho\ne0$. If $E_\epsilon<\infty$, then for $\sigma>\epsilon$ the Cauchy and Schwarz inequality gives $\int|D|e^{-\sigma u}du\le E_\epsilon^{1/2}/\sqrt{2(\sigma-\epsilon)}$, and with powers of $u$ the transform is holomorphic on $\Re z>\epsilon$, so no zero has $\Re\rho>\frac12+\epsilon$; for every $\epsilon$ this gives RH by the functional equation. Conversely, under RH von Koch's bound $\psi_{\mathrm C}(y)-y=O(\sqrt y\log^2y)$ gives $D(u)=O((1+u)^2)$ and $E_\epsilon<\infty$ for every $\epsilon>0$. For the limit, under RH $D(u)=-2\Re\sum_{\gamma>0}m_\gamma e^{i\gamma u}/(\frac12+i\gamma)+o(1)$ in the $B^2$ sense, and averaging the square of a finite sum over distinct frequencies leaves the squared coefficients, which gives the formula (Cram\'er \cite{Cramer} for the case of simple zeros; the multiplicity sensitive form is in Akbary, Ng and Shahabi, arXiv 1306.1657, unchecked). % UNCHECKED: Akbary, Ng and Shahabi, Theorem 1.14 and Corollary 1.15(i), taken from the GPT review of the source; the paper was not read for this book.
The unconditional identity $\sum_\rho1/(\rho(1-\rho))=2+\gamma_E-\log4\pi$, with multiplicity, comes from the logarithmic derivative of the Hadamard product at $0$ and at $1$.
\end{proof}

\begin{resultpar}\label{wt:res:energy}
Table~\ref{wt:tab:energy} checks the energy constant from the primes up to $2\times10^7$, integrating $D^2$ exactly on each interval between integers (computed). The last column is the sum under simple zeros, $0.046191$ (computed from the identity); finite windows can be consistent with it without proving the limit or the simplicity of the zeros.\end{resultpar}

\begin{table}[h]\centering\small
\fitwidth{\begin{tabular}{cccc}
\toprule
window in $u=\log y$ & mean of $D^2$, primes & Cram\'er model (one realisation) & $\sum_\rho|\rho|^{-2}$\\
\midrule
$[5,10]$ & 0.0509 & 5.19 & 0.0462\\
$[8,12]$ & 0.0442 & 3.43 & 0.0462\\
$[10,14]$ & 0.0465 & 2.93 & 0.0462\\
$[12,16.8]$ & 0.0468 & 5.14 & 0.0462\\
\bottomrule
\end{tabular}}
\caption{Energy of the normalised prime error on the log scale (computed). The source printed $0.0506$ for the first window and $7.54$, $9.78$, $9.14$, $7.15$ for a Cram\'er realisation; the Cram\'er values are heavy tailed and one realisation does not support a growth rate.}\label{wt:tab:energy}
\end{table}

This settles the variance question raised in Section~\ref{wt:sec:rigidity} in three parts.
\begin{enumerate}
\item The short interval variance of Goldston and Montgomery is a statement under RH, so it cannot be used by itself to prove RH.
\item A bound on the energy alone cannot prove $W(t)\ge0$ at a given $t$ when $\Sigma_\infty(t)<0$. If $E_\epsilon\le M^2$, weighted Cauchy and Schwarz gives $|\int Dk_t|\le MK_{t,\epsilon}$ with $K_{t,\epsilon}^2=\int_0^\infty k_t^2e^{2\epsilon u}du$, so $W(t)\in[\Sigma_\infty-MK_{t,\epsilon},\Sigma_\infty+MK_{t,\epsilon}]$. A small enough norm bound could certify a negative sign, and the limitation is exact for norm information: the two functions $D_\pm=\pm Me^{2\epsilon u}k_t/K_{t,\epsilon}$ have the same discounted energy and opposite extreme pairings. The norm depends on $t$ and is not uniformly of order one.
\item The global energy of $D$ is equivalent to RH (Proposition~\ref{wt:prop:energy}).
\end{enumerate}
Under RH the energy of $D$ per unit of $u$ is the sum of Proposition~\ref{wt:prop:energy}. The Cram\'er model has finite discounted energy too, almost surely for every $\epsilon>0$ at once: with $\Psi_C(y)=\sum_{3\le n\le y}B_n\log n$ and $B_n$ Bernoulli with parameter $1/\log n$, one has $\E\Psi_C(y)=y+O(1)$ and $\operatorname{Var}\Psi_C(y)=y\log y-2y+O(\log y)$, so $\E D_C(u)^2=O(1+u)$ and Tonelli gives $\E\int D_C^2e^{-2\epsilon u}du<\infty$. What it lacks is a bounded mean square per unit of $u$, which grows like $u$ in expectation. That this reflects the Euler product is a heuristic; the energy equivalence for the primes uses the zeta specific transform and is not a characterisation of Euler products.

\subsection{The reciprocal Xi law from primes}\label{wt:sec:cumulants}

Under RH the reciprocal $\Xi$ law $T=\sum_jE_j/\lambda_j$ has L\'evy measure $W(x)\,dx/x$ and Thorin measure $U=\sum_j\delta_{\lambda_j}$, so its cumulants are $\kappa_j=\int_0^\infty x^{j-1}W(x)\,dx=(j-1)!\,s_j$ with $s_j=\sum_k\lambda_k^{-j}=\sum_{\gamma>0}\gamma^{-2j}$. The integral identity $\int x^{j-1}W\,dx=(j-1)!\,s_j$ needs no RH: $\Re\lambda_j>0$ follows from $|\Im\gamma|<\frac12$ and the verification of RH to a finite height \cite{PT21}, and Fubini applies since $\sum_\gamma\int x^{j-1}|e^{-\gamma^2x}|dx=(j-1)!\sum(\Re\gamma^2)^{-j}<\infty$. So each $s_j$ is computable from $\Gamma$ and the primes alone. For a finite $X$, $\int_0^Xg_t(u)\,dt=\sqrt{X/\pi}\,e^{-u^2/4X}-\frac u2\operatorname{erfc}(u/2\sqrt X)$ for $u\ge0$, and the archimedean part integrates in closed form against $e^{t/4}$ and $g_t(0)$, leaving one integral against $\Re\psi(\frac14+\frac{ir}2)$. The source cut at $X=0.25$ and kept the prime powers up to $4000$; at 35 digits it printed the following (unchecked, not recomputed).
\begin{center}\small
\fitwidth{\begin{tabular}{ccc}
\toprule
& $\Gamma$ and primes $\le4000$ & from $\xi$ at $\frac12$\\
\midrule
$s_1=\sum\gamma^{-2}$ & 0.02310499311541897078893279 & 0.02310499311541897078893381\\
$s_2=\sum\gamma^{-4}$ & 0.00003717259928526968616461 & 0.00003717259928526968616487\\
\bottomrule
\end{tabular}}
\end{center}
The right column uses $\xi(\frac12+z)=\xi(\frac12)\prod(1+z^2/\gamma^2)$, so $s_1=\xi''(\frac12)/2\xi(\frac12)$ and $s_2=s_1^2-\xi^{(4)}(\frac12)/12\xi(\frac12)$; these are analytic coefficients, and calling them the mean and variance of $T$ is conditional on RH. The columns differ by $1.02\times10^{-24}$ and $2.6\times10^{-25}$, which are the omitted heat tails beyond $X=\frac14$: the first zero alone gives $\int_X^\infty W\ge e^{-\gamma_1^2X}/\gamma_1^2=1.0173\times10^{-24}$, so the cut is not below $10^{-24}$ as the source stated, and for $j=2$ the first zero gives $e^{-\gamma_1^2X}(X/\gamma_1^2+1/\gamma_1^4)=2.594\times10^{-25}$. The relative differences are $4\times10^{-23}$ and $7\times10^{-21}$, so the agreement is to 23 and 20 significant digits, not 24 for both. The omitted prime tail, bounded by the partial summation of Proposition~\ref{wt:prop:cert} applied to $G_X(u)=\int_0^Xg_t(u)\,dt$ (nonnegative and decreasing in $u$), is about $2\times10^{-31}$ at $N=4000$ (unchecked) and is far smaller.

This is a consistency check and not evidence for RH. What RH adds is that the atoms $\lambda_j^{-1}$ are real and positive, which is a moment problem. Put $m_j=s_{j+1}=\Gamma(j+1)^{-1}\int_0^\infty x^jW\,dx=\kappa_{j+1}/j!$, the normalised cumulants, and $\nu=\sum_j\lambda_j^{-1}\delta_{\lambda_j^{-1}}$, so that $m_j=\int t^j\,\nu(dt)$ and $|m_j|\le SR^j$ with $S=\sum|\lambda_j|^{-1}$ and $R=\sup|\lambda_j|^{-1}$.

\begin{proposition}[Hankel form of RH]\label{wt:prop:hankel}
RH holds if and only if the Hankel matrices $H_d=(m_{i+j})_{0\le i,j\le d}$ and $H_d^{(1)}=(m_{i+j+1})_{0\le i,j\le d}$ are positive semidefinite for every $d\ge0$.
\end{proposition}

\begin{proof}
Under RH $\nu$ is a positive measure and the two quadratic forms are $\int|p|^2d\nu$ and $\int t|p|^2d\nu$. Conversely, Stieltjes' theorem gives a positive measure $\mu$ on $[0,\infty)$ with these moments; from $b^{2k}\mu([b,\infty))\le m_{2k}\le SR^{2k}$ it is supported on $[0,R]$. Near $z=0$, $M(z)=\int(1-zt)^{-1}\mu(dt)=\sum_jm_jz^j=\sum_k1/(\lambda_k-z)$. The integral is holomorphic off $[1/R,\infty)$, while the spectral series is meromorphic, converging normally away from its poles, each of which has residue minus its positive integer multiplicity, so no cancellation occurs. Analytic continuation forces every $\lambda_k$ to be positive, and then $\Re\gamma_k>0$ with $\gamma_k^2>0$ gives $\gamma_k$ real, which is RH.
\end{proof}

Semidefiniteness means all quadratic forms, equivalently all principal minors; some leading determinants do not replace it. Each entry is a finite combination of integrals of $\mathcal A-\mathcal P$ against powers of $t$, convergent at the Gaussian rate in the primes. This is a criterion of the type of the Keiper and Li criteria and of Theorem~\ref{thm:ch12:hankel}, which has Hankel matrices at one heat time; here the entries are moments of the Thorin side. The extreme value structure of $W$ near $0$ (Section~\ref{wt:sec:gumbel}) makes the integrals converge rapidly.

\section{The sine squared lemma and a single heat time}\label{wt:sec:sin2}

The Thorin condition of Theorem~\ref{wt:thm:main} asks for every derivative of $W$ at every $t$. The sine squared lemma \eqref{eq:sin2}, in its L\'evy and Khintchine form $1-\cos\theta=2\sin^2(\theta/2)$ (valid for every complex $\theta$), trades that ladder for a one parameter family at a single heat time. The tent version used by Yoshida's window is Theorem~\ref{yo:thm:D}.

\subsection{The shifted functional}

For $t>0$ and $a\in\R$ put
\[
f_t(a)=\sum_je^{-\lambda_jt}\cos(a\gamma_j),\qquad V(t,a)=W(t)-f_t(a)=\sum_j2\sin^2\!\big(\tfrac{a\gamma_j}2\big)e^{-\lambda_jt},
\]
with the indexing of the introduction ($\Re\gamma_j>0$, multiplicities counted). For every complex $z$, $1-\cos az=\frac12(1-e^{iaz})(1-e^{-iaz})$. With $F(z)=(1-e^{iaz})e^{-tz^2/2}/\sqrt2$ and $F^*(z)=\overline{F(\bar z)}$, $F(z)F^*(z)=(1-\cos az)e^{-tz^2}$, which is $|F(z)|^2$ on the real axis. So $V$ is the Weil functional of the even part of an admissible test function $FF^*$ (it decays like a Gaussian on every strip). At a nonreal zero the summand need not be nonnegative (at $\gamma=ib$ the weight $2\sin^2(a\gamma/2)$ equals $-2\sinh^2(ab/2)<0$), so termwise nonnegativity holds only under RH. The explicit formula \eqref{wt:eq:ef} with $h(r)=e^{-tr^2}\cos(ar)$, for which $g(u)=\frac12\big(g_t(u-a)+g_t(u+a)\big)$, gives
\begin{multline}\label{wt:eq:fxa}
f_t(a)=e^{t/4}\cosh\tfrac a2-\tfrac12g_t(a)\log\pi+\frac1{4\pi}\int_\R e^{-tr^2}\cos(ar)\,\Re\psi\big(\tfrac14+\tfrac{ir}2\big)\,dr\\-\frac12\sum_{n\ge2}\frac{\Lambda(n)}{\sqrt n}\big(g_t(\log n-a)+g_t(\log n+a)\big).
\end{multline}
At $a=0$ this is $W=\mathcal A-\mathcal P$. For $a>0$ the shift moves each prime power atom from $\log n$ to $\log n-a$, so the primes near $e^a$ are brought to the centre of the Gaussian, while the pole term grows like $\frac12e^{t/4}e^{a/2}$ and must be cancelled by them. That cancellation is the prime number theorem near $y=e^a$, read through a Gaussian of width $\sqrt{2t}$ in $\log y$.

\begin{theorem}[One heat time suffices]\label{wt:thm:sin2}
Fix $t>0$. RH holds if and only if $V(t,a)\ge0$ for every $a\ge0$.
\end{theorem}

\begin{proof}
Under RH every $\gamma_j$ is real and each term of $V$ is nonnegative. Conversely let $V(t,\cdot)\ge0$ and suppose a zero is off the line, so that some $\gamma_j=\tau_j+i\beta$ has $\beta\ne0$; its conjugate occurs as well, so we may take $\beta>0$. Since $|\Im\gamma|<\frac12$, $|\cos(a\gamma)|\le e^{a/2}$ for $a\ge0$, and $\sum_j|e^{-\gamma_j^2t}|<\infty$ by the zero counting bound, the Laplace transform
\[
L(s)=\int_0^\infty e^{-sa}V(t,a)\,da=\frac{W(t)}s-\sum_je^{-\gamma_j^2t}\frac s{s^2+\gamma_j^2}
\]
holds for $\Re s>\frac12$. On each compact set avoiding the candidate poles the series converges normally (the summands are $O(e^{-\tau^2t}/\tau^2)$), so the right side is meromorphic on $\C$, with poles at $s=\pm i\gamma_j$, and a possible pole at $s=0$. The residue at a pole is $-\frac12$ times the multiplicity times $e^{-\gamma^2t}$, which is nonzero, and distinct ordinates with the same pole coincide, so poles do not cancel. The point $s=-i\gamma_j=\beta-i\tau_j$ is a genuine pole of real part $\beta>0$, so the abscissa of convergence $\sigma_c$ of the transform is at least $\beta>0$. Because $V\ge0$, Landau's theorem says that $L$ is singular at the real point $\sigma_c$: if $L$ were analytic there, choose $s_0=\sigma_c+d$ with $d>0$ small, so that the Taylor disk at $s_0$ has radius larger than $d$; its coefficients $(-1)^kL^{(k)}(s_0)/k!=\int a^ke^{-s_0a}V\,da/k!$ are nonnegative, and Tonelli sums the series at a real point $s_0-h$ with $h>d$ inside the disk to $\int e^{-(s_0-h)a}V\,da<\infty$, contradicting the definition of $\sigma_c$. The only real candidates for singularities of the right side are $s=0$ and the points $s=\pm i\gamma_j$ with $\gamma_j$ purely imaginary. The latter cannot occur: $\Re\gamma_j>0$ by the indexing, which is legitimate because $\xi$ has no real zeros ($\xi(s)>0$ for real $s$: for $0<s<1$ the alternating series gives $\eta(s)>0$, hence $\zeta(s)<0$ while $s(s-1)<0$, for $s>1$ the factors are positive, for $s<0$ use the functional equation, and $\xi(0)=\xi(1)=\frac12$). And $\sigma_c\ge\beta>0$ excludes $s=0$. This contradicts the existence of the real singularity, so every $\gamma_j$ is real.
\end{proof}

Where complete monotonicity used every order at every $t$, Theorem~\ref{wt:thm:sin2} uses one $t$ and every shift $a$; the shift plays the role of the order. By \eqref{wt:eq:smooth}, $f_t(a)$ is the Gaussian smoothing at scale $\sqrt t$ of the prime error $D$ near $u=a$, plus explicit archimedean terms, so the theorem is a one sided bound on a smoothed $D$ that is uniform in $a$. That is the quantity between pointwise and energy bounds left open in Section~\ref{wt:sec:round2}. It is close in spirit to classical facts about smoothed sums over zeros and is presented as a reformulation.

\subsection{Numerical check}

\begin{resultpar}\label{wt:res:sin2}
The source took $t=0.02$, the first 120 zeros and the prime powers up to 5000, at 30 digits (unchecked here; an independent calculation reproduced the six values to the printed digits from the first eight ordinates).\end{resultpar}
\begin{center}\small
\fitwidth{\begin{tabular}{ccc}
\toprule
$a$ & $V(t,a)$ from zeros & $V(t,a)$ from $\Gamma$ and primes, \eqref{wt:eq:fxa}\\
\midrule
0.3 & 0.0267570868155 & 0.0267570868155\\
$\log2$ & 0.0357321878430 & 0.0357321878430\\
$\log3$ & 0.0367065917181 & 0.0367065917181\\
1.5 & 0.0313514051449 & 0.0313514051449\\
$\log8$ & 0.0264484602913 & 0.0264484602913\\
3.0 & 0.0185315145626 & 0.0185315145626\\
\bottomrule
\end{tabular}}
\end{center}
The source reported a minimum of $V/W$ of $4.5\times10^{-4}$ on a grid of 1600 points in $0<a\le8$ (unreproduced; on the grid $a=8k/1600$ an independent calculation gives about $5.6\times10^{-4}$ at $a=0.89$, and other grids give other values, none a continuous minimum). At this $t$ the lowest zero carries $99.2\%$ of $W$ ($0.99198$; the source's $99.9\%$ is not possible, because the second zero alone brings the share below $0.99218$). The local minima sit near $a\approx2\pi k/\gamma_1=0.4445k$ and are deep for even $k$ (ratio about $4.9\times10^{-4}$ at $k=2$ and $4.2\times10^{-4}$ at $k=4$, unchecked) and shallow for odd $k$ (about $0.016$, unchecked), since $2\gamma_2/\gamma_1=2.975$ is close to an integer. For smaller $t$ at a fixed positive shift the prime power resonance $a_0=\log m$ gives a peak, not a dip: $V(t,a_0)=W(t)+\Lambda(m)/(4\sqrt{\pi mt})+O(1)$ as $t\downarrow0$.

To see the criterion fail, the source moved $\gamma_1$ off the line, replacing it by the pair $\gamma_1\pm i\beta$ and keeping the other zeros. The pair contributes $2\Re[c(1-\cos(a(\gamma_1+i\beta)))]$ with $c=e^{-t(\gamma_1+i\beta)^2}$, whose leading oscillation at large $a$ is $-e^{-t(\gamma_1^2-\beta^2)}e^{\beta a}\cos(a\gamma_1+2t\gamma_1\beta)$. For $\beta=0.05$ the minimum of $V$ over $0<a\le200$ is of the order of $-4\times10^2$ (the source printed $-391$ on a coarse sample; the single shift $a\approx199.588$ gives $-396.8$ in an independent calculation), and for $\beta=0.2$ it is about $-4\times10^{15}$, both near $a=200$ (unchecked here). This is an experiment on one finite model, not a proof of a general domination.

\subsection{What it costs}

The prime side of $V(t,a)$ is governed by the prime powers near $e^a$ in $\log n$. Write $V=V_N-\sum_{m>N}\Lambda(m)m^{-1/2}H_{t,a}(\log m)$ with $H_{t,a}(u)=g_t(u)-\frac12[g_t(u-a)+g_t(u+a)]=g_t(u)\big[1-e^{-a^2/4t}\cosh(au/2t)\big]$ and $V_N$ the same expression with the sum truncated at $N$ and the archimedean part written out. For $a>0$, $H_{t,a}(u)\le0$ for $u\ge u_0(t,a)=(2t/a)\operatorname{arcosh}(e^{a^2/4t})$, so $\log N\ge u_0$ implies the rigorous one sided bound $V\ge V_N$; for $0\le a\le A$ the condition $\log N\ge A+\sqrt{2t}$ suffices, by convexity of $g_t$ on $[\sqrt{2t},\infty)$. At $a=0$, $V=0$ exactly, so a positive uniform margin is impossible, and the quotient $V(t,a)/a^2$, whose value at $0$ is $-W'(t)/2$, is what a certificate must enclose. This proved extension of the certificate argument needs the tail hypotheses checked as in Proposition~\ref{wt:prop:cert}; where the sign condition is unavailable an absolute tail majorant with partial summation is used. The number of primes needed is not simply "up to $e^{a+O(\sqrt t)}$": to keep a fixed small absolute tail $\eta$ at large $a$, completing the square gives $\int_L^\infty e^{u/2}g_t(u-a)\,du=\frac12e^{a/2+t/4}\operatorname{erfc}\big((L-a-t)/2\sqrt t\big)$ and so $L=\log N=a+t+O(\sqrt{t(a+\log(1/\eta))})$ in the bound of that type, and a one sided certificate that drops a helpful signed tail can use a different cutoff. For large $a$ the condition is square root control of $\psi_{\mathrm C}$ near $y=e^a$ in a smoothed one sided sense, which is not literally von Koch's pointwise estimate. The obstruction moves from large $t$ to large $a$ but is not removed; the criterion is local in $a$, and a failure of RH must show up as a negative value of a single explicit function of one variable.
